\section{Order, incidence, and extremal set systems}
\label{sec:order-structure}

The survey asks for the largest antichain and the largest intersecting uniform
family on an \(n\)-element set. The answers use two different ways to read
a permutation: take its prefixes, or wrap it around a circle. We develop
the order-theoretic results first. Throughout, a subset of \([n]=\{0,\ldots,n-1\}\)
is a binary mask, a family is a mask of masks, and a poset on \([m]\) is
a reflexive, antisymmetric, transitive Boolean table. An integer can be
coded by a sign and an absolute value, with one code for zero.

\subsection{Chains, antichains, and finite lattices}

\begin{theorem}[Linear extensions and Mirsky's theorem]\label{result35}
A finite poset has a computable linear extension. Its maximum chain size
equals the minimum number of antichains partitioning it. Both extrema
come with witnesses.
\end{theorem}
\begin{proof}
For \(x\), let \(d(x)\) be the number of its weak predecessors. If \(x<y\),
the predecessor mask of \(x\) is properly included in that of \(y\), with
\(y\) witnessing properness. Thus \(d(x)<d(y)\). Sort by \((d(x),x)\)
lexicographically to get a linear extension. Read it in order and write
\[
 h(x)=1+\max(\{0\}\cup\{h(z):z<x\}).
\]
Induction on the processed prefix proves simultaneously that every chain
ending at a processed \(x\) has at most \(h(x)\) elements and that a chain
of exactly this size exists. For the upper bound, remove the endpoint
of a nonsingleton chain. For the lower bound append \(x\) to the recorded
chain at the first maximizing predecessor, using the singleton when the
predecessor set is empty. Strict comparison increases \(h\), hence every
nonempty fiber is an antichain. A longest chain meets distinct fibers
and needs a different part for each element in any antichain partition.
For the empty poset both numbers are zero. Chains are arrays of at most
\(m\) distinct labels; testing the upper-bound invariant amounts to a
bounded scan of padded arrays over \([m+1]\).
\end{proof}

\begin{theorem}[Dilworth's theorem]\label{result36}
The largest size of an antichain is the least number of chains partitioning
a finite poset. A partition and an antichain of the same size can be constructed.
\end{theorem}
\begin{proof}
On two copies of \([m]\) put the bipartite edges \(x_Ly_R\) with \(x<y\).
The matching and cover algorithm of Theorem~\ref{thm:matching-cover}
produces a matching and a vertex cover of the same size \(e\).
Treat each matched edge as an arrow \(x\to y\) on the original labels.
Each vertex has at most one incoming and one outgoing arrow. Each arrow
strictly increases \(d\), so following predecessors reaches a start
within \(m\) steps, and following successors from all starts partitions
the vertices into chains. There are \(m-e\) starts because exactly \(e\)
vertices have a predecessor. Keep labels \(x\) for which neither copy
belongs to the cover. At most \(e\) distinct labels were discarded.
Two kept labels cannot satisfy \(x<y\), since their edge would miss the
cover. This gives an antichain of size at least \(m-e\), and no antichain
can meet one of the constructed chains twice. These two witnesses
also certify both optimality claims.
\end{proof}

\begin{corollary}[Erd\H{o}s--Szekeres]\label{result37}
A list of \(rs+1\) pairwise distinct elements of a coded decidable linear order has a strictly increasing
subsequence of length \(r+1\) or a strictly decreasing one of length \(s+1\).
\end{corollary}
\begin{proof}
For each position \(i\) record the longest increasing and decreasing
subsequences ending there using the preceding recurrence over earlier
positions with smaller or larger value, respectively. If \(i<j\), either
the first length increases strictly from \(i\) to \(j\), or the second
does. All pairs of lengths differ. If all first lengths are at most \(r\)
and all second lengths at most \(s\), there are only \(rs\) pairs of
positive lengths. The pigeonhole count rules this out; scan for the first
excess and recover a subsequence from its predecessor pointers.
\end{proof}

\begin{theorem}[Finite distributive lattice representation]\label{result38}
A nonempty finite distributive lattice \(L\) is isomorphic to the lattice
of downward closed subsets of its poset \(J\) of nonzero join-irreducibles.
The isomorphism, its inverse, and \(J\) are primitive recursive in the
supplied join and meet tables.
\end{theorem}
\begin{proof}
The order is \(u\le v\) when \(u\vee v=v\); the bottom \(0_L\)
is the meet of all elements. A nonzero \(j\) belongs to \(J\) when
\(j=u\vee v\) always implies \(j=u\) or \(j=v\), a test over table indices.
Write \(I_x=\{j\in J:j\le x\}\). Every \(x\) is the join of \(I_x\)
(the empty join is \(0_L\)). Prove this in the linear extension:
for \(0_L\) and for \(x\in J\) it follows immediately. Otherwise a bounded
scan supplies \(x=u\vee v\) with \(u,v<x\); insert the two earlier
decompositions. Their terms all belong to \(I_x\), and all elements of
\(I_x\) lie below \(x\).
Distributivity makes each \(j\in J\) join-prime: if \(j\le u\vee v\)
then \(j=(j\wedge u)\vee(j\wedge v)\), so one summand equals \(j\).
Iteration gives the same conclusion for a nonempty finite join; a
nonzero \(j\) cannot lie below an empty join. For a downward closed
\(D\subseteq J\), put \(x=\bigvee D\). Then \(D\subseteq I_x\) and
each \(j\in I_x\) is below a member of \(D\), hence belongs to \(D\).
We have both inverse identities. Finally \(I_{u\wedge v}=I_u\cap I_v\)
by the meet property, and \(I_{u\vee v}=I_u\cup I_v\) by join-primality.
\end{proof}

\subsection{Finite incidence inversion}

\begin{theorem}[M\"obius inversion and product law]\label{result39}
In a coded finite poset \(P\), form the signed total \(\mu_P(x,y)\)
of the strict chains \(x=x_0<\cdots<x_j=y\), weighted by \((-1)^j\).
The singleton chain has weight \(1\). For integer tables \(f,g\),
\[
 g(y)=\sum_{x\le y} f(x)\quad\hbox{for every }y
 \quad\Longleftrightarrow\quad
 f(y)=\sum_{x\le y}\mu_P(x,y)g(x)\quad\hbox{for every }y.
\]
Also \(\mu_{P\times Q}((x,u),(y,v))=\mu_P(x,y)\mu_Q(u,v)\).
For \(A\subseteq B\) in the Boolean lattice this gives
\(\mu(A,B)=(-1)^{|B\setminus A|}\).
\end{theorem}
\begin{proof}
For \(x<y\), sum the chains contributing to
\(\sum_{x\le z\le y}\mu_P(x,z)\). If their last point is \(y\), delete
it; otherwise append \(y\). This involution reverses the sign. For the
sum of \(\mu_P(z,y)\), insert or delete \(x\) at the beginning instead.
For \(x=y\) both sums are \(1\). For a function supported on comparable
pairs, define convolution by
\((a*b)(x,y)=\sum_{x\le z\le y}a(x,z)b(z,y)\).
The interval indicator \(\zeta(x,y)=[x\le y]\) therefore has a two-sided
inverse \(\mu_P\). Associativity follows by enumerating the same pairs
\(x\le z\le w\le y\) in either order, and the diagonal indicator acts
as identity. Expanding each nested sum in the displayed equivalence
gives the two inverse identities just proved.
On a product interval the intermediate coordinates vary independently;
hence convolution of separate products factors into the two coordinate
convolutions. The proposed product of M\"obius values thus inverts the
product zeta; a two-sided inverse is unique by associativity.
The two-point chain has diagonal entries \(1\) and strict entry \(-1\);
iterating the product law for \(n\) copies proves the subset formula
(including the zero-fold product). The chain sum on \([m]\) can be computed
by sorting the vertices of each of its \(2^m\) masks in the fixed linear
extension and checking endpoints and strict comparisons. Its absolute
value is therefore at most \(2^m\), a bound that also applies to
intermediate signed totals.
\end{proof}

\subsection{Maximal chains and cyclic intervals}

\begin{theorem}[LYM, Sperner, and the general chain bound]\label{result40}
Let \(\mathcal F\) be a family of subsets of \([n]\) with at most \(r\)
members in any strict chain, \(0\le r\le n+1\). Put \(a_j=\#\{A\in
\mathcal F:|A|=j\}\) and \(c_j=\binom nj\). Then
\[
 \sum_{j=0}^n a_j\,j!(n-j)!\ \le\ r\,n!.
\]
The maximum size of such a family is the sum of the \(r\) largest
\(c_j\). In particular the maximum antichain size is
\(\binom n{\lfloor n/2\rfloor}\), and every maximizing antichain is
one full level of that size.
\end{theorem}
\begin{proof}
Each linear ordering of \([n]\) specifies one chain of all its prefixes,
including the empty prefix. A particular \(j\)-set appears in
\(j!(n-j)!\) orders: order it first and its complement last. Count
the pairs (order, selected prefix) by their two coordinates to obtain
the inequality. Dividing formally by \(n!\) writes it as
\(\sum_j a_j/c_j\le r\); all denominators divide \(n!\).
For \(r=0\) the family is empty. For \(r>0\), let \(T\) be \(r\) levels
of largest size and \(c=\min_{j\in T}c_j\). Then
\[
 \sum_{j\notin T}a_j
 \le c\sum_{j\notin T} a_j/c_j
 \le c\sum_{j\in T}(1-a_j/c_j)
 \le \sum_{j\in T}(c_j-a_j).
\]
On a chain there is at most one member of each rank, so the union of
these \(r\) levels attains the bound. Comparing \(c_{j+1}(j+1)=
c_j(n-j)\) identifies the middle maximum (two equal levels for odd \(n\)).

For an antichain of maximum size \(M=\max c_j\), the inequalities
\(\sum_j a_j/M\le\sum_j a_j/c_j\le1\) force every selected level
to have size \(M\), and every prefix chain to contain a selected set.
For even \(n\) there is one candidate level. For odd \(n=2h+1\),
join an \(h\)-set to an \((h+1)\)-set containing it. Each such pair
extends to a prefix chain, so exactly one endpoint is selected. This
bipartite graph is connected: two \(h\)-sets can be reached from each
other by exchanging one unwanted element for one wanted element at a
time, walking through their union at each exchange. Every upper set has
a lower neighbor, including when \(h=0\). Hence the choice at one lower
vertex forces the choice on the whole graph. This also proves the
classification for \(n=1\).
\end{proof}

\begin{theorem}[Erd\H{o}s--Ko--Rado]\label{result41}
For \(1\le k\) and \(2k\le n\), the largest pairwise intersecting
family \(\mathcal F\) of \(k\)-subsets of \([n]\) has size
\(\binom{n-1}{k-1}\). For \(k\le n<2k\), its maximum is \(\binom nk\).
\end{theorem}
\begin{proof}
Place the positions of a permutation on a circle. At most \(k\) of its
length-\(k\) arcs can belong to \(\mathcal F\). To see this, handle
the empty case and otherwise rotate the labels of positions so that a
selected arc starts at \(0\). Each other selected start lies between
\(-(k-1)\) and \(k-1\) modulo \(n\), since its arc meets the first one.
Pair start \(-i\) with \(k-i\) for \(1\le i<k\). Their arcs are
disjoint: they occupy \(2k\) consecutive positions without repetition.
Thus at most one of each of \(k-1\) pairs can be selected.
For each \(A\in\mathcal F\) and each of the \(n\) possible starts,
exactly \(k!(n-k)!\) permutations display \(A\) on that arc. Counting
(permutation, start) gives
\(|\mathcal F|n k!(n-k)!\le k\,n!\), the claimed bound after
cancellation of positive factors. All \(k\)-sets through one fixed
label meet and attain it. When \(2k>n\), two disjoint \(k\)-sets
would have \(2k\) distinct members of \([n]\), so every pair meets.
\end{proof}

\section{Designs, finite planes, and error-correcting codes}
\label{sec:incidence-structure}

\subsection{Parameter arithmetic and Fisher's inequality}

A \(t\)-\((v,k,\lambda)\) design is an indexed list of \(k\)-subsets
of \([v]\) in which every \(t\)-subset occurs in exactly \(\lambda\)
blocks. We use \(0\le t\le k\le v\). Indexing makes multiplicities
explicit.

\begin{theorem}[Descent, incidence counts, and derived designs]\label{result42}
Every \(s\)-subset, \(0\le s\le t\), belongs to the same number
\(\lambda_s\) of blocks, determined by
\[
 \lambda_s\binom{k-s}{t-s}=\lambda\binom{v-s}{t-s}.
\]
In a \(2\)-design with \(b\) blocks and replication \(r=\lambda_1\),
\(bk=vr\) and \(r(k-1)=\lambda(v-1)\).
For \(t\ge 1\), the blocks through a specified point, with that point
deleted, form a \((t-1)\)-\((v-1,k-1,\lambda)\) design.
The blocks missing the point form a \((t-1)\)-design with block size
\(k\) and parameter \(\lambda_{t-1}-\lambda\) when \(k<v\).
\end{theorem}
\begin{proof}
Fix an \(s\)-set \(S\) and count the pairs \((T,B)\) with
\(S\subseteq T\subseteq B\) and \(|T|=t\). Choosing the extra
\(t-s\) points before choosing a block gives the right side; choosing
a block through \(S\) first gives the left with the actual number
through \(S\) as multiplier. The other factor is positive, so these
actual numbers coincide by cancellation. Count pairs (point, block
through it) to get \(bk=vr\); descent at \(s=1,t=2\) gives the second
identity. For any \((t-1)\)-set avoiding the specified point, adjoining
that point produces a \(t\)-set in exactly \(\lambda\) blocks. Of the
\(\lambda_{t-1}\) blocks containing the smaller set, these are exactly
the ones removed in the residual design and retained in the derived design.
\end{proof}

\begin{theorem}[Fisher's inequality]\label{result43}
A \(2\)-\((v,k,\lambda)\) design with \(2\le k<v\) and \(\lambda>0\)
has \(b\ge v\).
\end{theorem}
\begin{proof}
The incidence identities give \(r>\lambda\). For integral coordinates
\(z_0,\ldots,z_{v-1}\), expansion and the two incidence multiplicities give
\[
 \sum_B\left(\sum_{i\in B}z_i\right)^2
  =(r-\lambda)\sum_i z_i^2+\lambda\left(\sum_i z_i\right)^2.
\]
Indeed a square \(z_i^2\) has total coefficient \(r\), and each product for
\(i<j\) has coefficient \(2\lambda\). If all block sums vanish, scanning
the squares on the right shows that all coordinates vanish. Consequently
the block sums map is injective on any box of integer vectors, by taking
the difference of two vectors.

If \(b<v\), choose \(M=(v+1)^b+1\) and restrict to coordinates in
\(\{0,\ldots,M-1\}\). The domain has \(M^v\) vectors. Each block sum
has fewer than \((v+1)M\) possible positive-or-zero values, and the
number of image vectors is at most
\(((v+1)M)^b<M^{b+1}\le M^v\).
This contradicts finite pigeonhole. Since \(b<v\) is decidable, \(v\le b\)
follows. Even the auxiliary box used here has a specified primitive
recursive size.
\end{proof}

\begin{corollary}[Necessary Steiner triple parameters]\label{result44}
If \(v\ge3\) points carry a family of triples through every pair exactly once,
then \(v\equiv1\) or \(3\pmod6\), each point is in \((v-1)/2\) triples,
and there are \(v(v-1)/6\) triples. A parallel class is a partition of all
points into triples. If the triples partition into parallel classes, then
\(v\equiv3\pmod6\) and the number of classes is \((v-1)/2\).
\end{corollary}
\begin{proof}
Take \(k=3,\lambda=1\). The equations \(2r=v-1\) and \(3b=vr\)
force \(v\) odd and \(3\mid v(v-1)\); checking residues modulo \(6\)
leaves \(1,3\). A parallel class requires \(3\mid v\), and uses a
given point exactly once.
\end{proof}

\subsection{An explicit solution to Kirkman's schoolgirl problem}

\begin{proposition}[Resolvable Steiner triple system on fifteen points]\label{result45}
Fifteen pupils can walk in five disjoint triples on each of seven days
so that each pair walks together exactly once.
\end{proposition}
\begin{proof}
Label them \(\infty,A_i,B_i\) with \(i\in\mathbb Z/7\mathbb Z\). Use
the following five rows on day \(0\), and add \(d\bmod7\) to every finite
subscript on day \(d\).
\[
\begin{array}{ccccc}
(\infty,A_0,B_0)&(A_1,A_2,A_4)&(A_3,B_1,B_5)&
(A_5,B_4,B_6)&(A_6,B_2,B_3).
\end{array}
\]
Each day is a partition by inspection of the subscripts. Translation of the
first row accounts for all pairs involving \(\infty\). Up to sign, the
differences in the all-\(A\) row are \(1,2,3\); the three \(B\)-pairs
in the last three rows have differences up to sign \(3,2,1\).
The mixed differences \(j-i\) from the displayed \(A_i,B_j\) pairs are
\(0;\ 5,2;\ 6,1;\ 3,4\), each residue once. For each mixed difference
its seven translates give precisely that class of pairs; for each
nonzero same-letter difference up to sign the seven translates do likewise,
since \(d\not\equiv-d\pmod7\) when \(d=1,2,3\). This covers every
unordered pair exactly once.
\end{proof}

\subsection{Affine and projective planes over a finite field table}

\begin{theorem}[Planes of order \(q\)]\label{result46}
Given a field table \(F\) on \(q\ge2\) elements, there is an affine plane
with \(q^2\) points and \(q(q+1)\) lines of size \(q\), arranged in
\(q+1\) parallel classes. Adjoining the points at infinity constructs a
projective plane with \(q^2+q+1\) points and as many lines, all of size
\(q+1\), with \(q+1\) lines through each point. In particular both
constructions are uniform in a prime \(p=q\).
\end{theorem}
\begin{proof}
On \(F^2\) use the \(q\) verticals \(x=a\) and the \(q^2\) lines
\(y=mx+c\). A vertical has one point for each \(y\), and a graph
one for each \(x\). Two distinct points with equal \(x\)-coordinate
lie on exactly their vertical. With unequal \(x\)-coordinates they determine
\(m=(y'-y)/(x'-x)\) and \(c=y-mx\). For a fixed slope the \(q\) graphs
partition the points, as do the \(q\) verticals. Lines of different
nonvertical slopes intersect at the unique solution of
\((m-m')x=c'-c\); a vertical meets each graph once. This proves all
affine incidence and parallel axioms. The points \((0,0),(0,1),(1,0)\)
supply a noncollinear triple.

Add a point at infinity for each slope including the vertical direction,
append it to every line of that class, and make all \(q+1\) new points
one new line. Two ordinary points retain their unique joining line.
An ordinary and an infinite point specify the unique line through
the former of that direction. Two infinite points determine the new
line. Distinct old lines in one class now meet at their infinite point,
distinct classes still meet at their original point, and the new line
meets an old line at its direction. There are \(q^2+q+1\) points,
\(q(q+1)+1\) lines, and the displayed incidence verifies the stated
valencies. The four corners \((0,0),(0,1),(1,0),(1,1)\) have no
three collinear, since any three include an entire vertical side and a
point outside it. This gives the nondegeneracy axiom. For \(p\) prime
use arithmetic modulo \(p\). For \(0<a<p\) the Euclidean algorithm
on \(p,a\) terminates within \(p+1\) divisions (positive remainders
strictly decrease); back-substitution gives \(up+va=1\), since the
last nonzero remainder divides the prime \(p\) and also \(a\).
Thus scanning residues for an inverse of \(a\) succeeds.
\end{proof}

\begin{corollary}\label{result47}
An affine plane from this construction is a resolvable
\(2\)-\((q^2,q,1)\) design. A projective plane from it is a symmetric
\(2\)-\((q^2+q+1,q+1,1)\) design (as many blocks as points).
At \(q=2\) it gives the seven-point, seven-line Fano plane.
\end{corollary}
\begin{proof}
Use the lines as blocks and apply uniqueness of the joining line.
For an explicit table at \(q=2\), let \(a,b,c,d\) denote respectively
\(00,01,10,11\); let \(v,h,s\) denote vertical, horizontal and
slope-one infinity. The seven lines are
\(abv,cdv,ach,bdh,ads,bcs,vhs\).
\end{proof}

\subsection{Metric bounds and the perfect binary Hamming code}

A \(q\)-ary word of length \(n\) is a list of \(n\) digits in \([q]\),
\(q\ge2\). Hamming distance counts positions with unequal digits. A code
is a mask of the \(q^n\) words. Write
\[
 V_q(n,t)=\sum_{i=0}^{\min(n,t)}\binom ni(q-1)^i .
\]

\begin{theorem}[Hamming, Singleton, and Gilbert bounds]\label{result48}
Let \(1\le d\le n+1\), and let \(C\) have pairwise distances at least
\(d\). For \(t=\lfloor(d-1)/2\rfloor\),
\[
 |C|V_q(n,t)\le q^n,\qquad |C|\le q^{\,n-d+1}.
\]
There is a computable code of pairwise distances at least \(d\) with
\(\lvert C\rvert V_q(n,d-1)\ge q^n\).
\end{theorem}
\begin{proof}
At distance \(i\) from a center, choose \(i\) changed coordinates and
one of \(q-1\) other digits in each, giving the stated ball volume.
Coordinatewise, a difference at the endpoints requires a difference on
one leg of any two-step path, proving the triangle inequality by addition.
Radius-\(t\) balls about codewords are disjoint because their common
point would force distance at most \(2t<d\). Counting all their points
gives the first bound. Delete a fixed \(d-1\) coordinates from every
codeword. Two equal resulting words would have differed only in those
coordinates; hence deletion is injective on \(C\), proving Singleton's
bound. For Gilbert's lower bound, visit the words in numerical order,
retaining a word if its distance to every retained word is at least \(d\).
Separation is preserved at each insertion. Each rejected word was within
\(d-1\) of an earlier retained word and remains covered. Thus balls of
radius \(d-1\) around the final code cover all \(q^n\) words. Counting
them with multiplicity gives the opposite inequality.
\end{proof}

\begin{theorem}[Perfect binary codes and their triples]\label{result49}
For \(r\ge2\) and \(n=2^r-1\), there is a binary linear code of length
\(n\), size \(2^{n-r}\), and minimum distance \(3\), whose radius-one
balls partition all words. Its words of weight three have supports forming
a Steiner triple system on \(n\) labels.
\end{theorem}
\begin{proof}
Index coordinates by the nonzero \(r\)-bit words. Let \(\sigma(w)\)
be the bitwise exclusive-or of labels whose coordinate in \(w\) is \(1\);
let \(C=\{w:\sigma(w)=0\}\). Exclusive-or is coordinatewise addition
modulo \(2\) and cancellation in pairs shows that \(\sigma\) preserves
addition. No one- or two-coordinate nonempty support has zero syndrome:
one label is nonzero, and two distinct labels have nonzero sum. Every
pair \(a,b\) of distinct nonzero labels yields exactly one third
label \(a\mathbin{\oplus}b\) completing a zero-syndrome triple, distinct
from \(0,a,b\). Pairs exist for \(r\ge2\). By linearity the minimum
distance therefore equals three, and the triples have the asserted
unique-pair property.

Given \(w\), leave it if \(\sigma(w)=0\), and otherwise flip the unique
coordinate labeled \(\sigma(w)\). Its resulting syndrome is zero, and a
single flip at any other coordinate would leave a nonzero syndrome. This
proves unique decoding at radius one. For an independent count of the code,
first assign arbitrary bits to all positions except the \(r\) coordinate
unit vectors. The syndrome of the partial assignment specifies the unique
bits to put in those \(r\) positions, one per set bit of that syndrome.
Hence \(|C|=2^{n-r}\), and indeed \((1+n)|C|=2^n\).
\end{proof}

\subsection{Lean definitions and induction obligations for these results}

The following specifications use bounds as arguments. They can be
implemented over natural-number tables using the coding layer of this
paper. The named lemmas are a dependency order for completing the Lean
proofs. The counting arguments use equality of interval sums under
explicit rearrangements. Divisions in combinatorial inequalities abbreviate
multiplication by positive common denominators.

\paragraph{Order and incidence.}
Use \(\mathrm{le}(P,x,y)\) for bit \(mx+y\) of a relation mask.
Define \(\mathrm{validPoset}(m,P)\) by bounded diagonal, pair and
triple scans. Define \(\mathrm{predCount}\) by summation and
\(\mathrm{extension}\) as selection of the unused least
\((\mathrm{predCount},x)\) for exactly \(m\) iterations.
Prove \texttt{predCount\_lt} using subset counting, then
\texttt{extension\_perm} and \texttt{extension\_sorted} by prefix
induction. Store the longest-chain lengths and first maximizing predecessors
in parallel tables. Induct on the extension prefix with table entries,
pair comparisons, and all candidate chains bounded by
\((m+1)^{m+1}\) using a padding digit. For Dilworth, construct the two
copies of the poset via \(x\mapsto x\) and \(x\mapsto m+x\); reuse
\texttt{matching-cover} (Theorem~\ref{thm:matching-cover}) and iterate the
matched successor \(m\) times with an end marker.
The lattice functions \texttt{joinIrred}, \texttt{idealOf} and
\texttt{joinMask} require scans of at most \(m^2, m,m\) indices for a
fixed element or mask. Prove the reconstruction identity first, the
join-prime test next, then the inverse identity by induction on a join
prefix. For M\"obius inversion, define \texttt{chainMask},
\texttt{signedChainCount}, \texttt{conv}; prove the two endpoint-toggle
cancellations, convolution associativity by summation over the same two
intermediate indices with endpoints fixed, and the factored double sum. This yields the inversion,
product, and Boolean formulas.

\paragraph{Extremal families.}
Define \texttt{prefixMask} and \texttt{arcMask} on permutation tables
of length \(n\) with digits below \(n\). For each mask of size \(k\), enumerate
the permutations containing it in a prescribed block by packing the two
separately ordered restrictions into the block and its complement, each
in increasing position order. Their unpacking maps are inverses by
digit inspection. Prove \texttt{prefix\_fiber} and
\texttt{arc\_fiber}; the latter regards the selected arc as the list
of positions from its start modulo \(n\). For LYM, sum the prefix
fiber bound and sort the \(n+1\) numerical rank sizes to select \(T\).
Test the chain restriction using all prefix permutations: every chain
can be extended to one by listing successive set differences in increasing
order and then the outside elements. For EKR, prove bounded
\texttt{pairedArcs\_disjoint} for \(1\le i<k\) by the \(2k\) consecutive
positions argument, and sum its indicator inequalities.

\paragraph{Designs and geometry.}
Define \texttt{blockAt}, \texttt{through} and \texttt{designOK} using
a list of \(b\) masks and a \(2^v\)-mask scan. The first reusable lemma
counts \((T,B)\) by transposing its bounded double sum. For Fisher's
bound define \texttt{blockSums} on \(v\) digits with base \(M\), prove
\texttt{sumSquares\_expand} by finite distributivity, then
\texttt{blockSums\_injective} by integer subtraction and positivity,
and compare the two box cardinalities. For Kirkman use \texttt{day d}
formed from the five displayed triples. The reduction to residues gives
seven-bit difference tests or, equivalently, \(105\) unordered-pair
tests on the full schedule. The supplied \texttt{lean/Kirkman.lean}
checks the latter. Its theorems \texttt{shape}, \texttt{attendance},
and \texttt{pairCoverage} quantify over \texttt{Fin 7} and \texttt{Fin 15};
Lean discharges all three by kernel reduction and reports no axiom dependencies. For field tables, encode points by \(qx+y\),
lines by \(qm+c\) followed by the \(q\) verticals. Append the \(q+1\)
infinite points and one new line. \texttt{joinPoints} and
\texttt{meetLines} branch on the types and equalities specified in the
proof. A scan of \(q\) labels computes an inverse of a nonzero field
element. Check existence and uniqueness by substitution in the table
laws, expressed by bounded scans over \(q^3\) triples.

\paragraph{Codes.}
Define \texttt{dist n q u v} as a sum of digit inequalities;
\texttt{syndrome r w} as a fold over the \(2^r-1\) nonzero labels;
\texttt{correct} flips the syndrome-labeled bit when the syndrome is
nonzero. Prove \texttt{dist\_triangle} by a sum of Boolean inequalities,
\texttt{ball\_card} by enumerating change masks and changed digits,
\texttt{syndrome\_flip} by cancellation, and the two-sided specification
of the encoder by comparing the \(r\) component bits. The Gilbert loop
uses \(q^n\) iterations; its invariant says that all retained words are
separated and that every processed word lies within \(d-1\) of some
retained word. Both assertions are bounded by the word range \(q^n\).
All induction assertions above are equalities or inequalities of
primitive recursive table checks, with codes supplied as parameters.

\paragraph{Sources.}
The survey's order, design, finite geometry, extremal, and coding sections
motivate these results. Its extremal questions and its Kirkman example
are answered explicitly. The incidence-algebra calculation follows
G.-C. Rota, \emph{On the foundations of combinatorial theory I}
(\href{https://doi.org/10.1007/BF00531932}{1964}).
The permutation and circle counts follow D. Lubell,
\emph{A short proof of Sperner's lemma}
(\href{https://doi.org/10.1016/S0021-9800(66)80035-2}{1966}),
and G. O. H. Katona, \emph{A simple proof of the Erd\H{o}s--Chao
Ko--Rado theorem}
(\href{https://doi.org/10.1016/0095-8956(72)90054-8}{1972}).
For the original binary code construction see R. W. Hamming,
\emph{Error detecting and error correcting codes}
(\href{https://doi.org/10.1002/j.1538-7305.1950.tb00463.x}{1950}).
